% Part: computability
% Chapter: recursive-functions
% Section: other-recursions

\documentclass[../../../include/open-logic-section]{subfiles}

\begin{document}

\olfileid{cmp}{rec}{ore}
\olsection{অন্যান্য পুনরাবৃত্তি}

যুগলায়ন ও অনুক্রমায়ন ব্যবহার করে আদিম পুনরাবৃত্তির আরও অভিনব এবং
উপযোগী রূপগুলিকে ন্যায্য প্রতিপন্ন করা যায়। যেমন, নিচের সংজ্ঞার মতো
একই সঙ্গে দুটি অপেক্ষক সংজ্ঞায়িত করা প্রায়ই উপযোগী:
\begin{align*}
h_0(\vec x, 0) & = f_0(\vec x) \\
h_1(\vec x, 0) & = f_1(\vec x) \\
h_0(\vec x, y+1) & = g_0(\vec x, y, h_0(\vec x, y), h_1(\vec x, y)) \\
h_1(\vec x, y+1) & = g_1(\vec x, y, h_0(\vec x, y), h_1(\vec x, y))
\end{align*}
এটি \emph{যুগপৎ পুনরাবৃত্তি}-র একটি নিদর্শন। অপেক্ষক সংজ্ঞায়নের আরেকটি
উপযোগী উপায় হলো $h(\vec x, y+1)$-এর মানকে
\emph{সমস্ত} মান $h(\vec x, 0)$, \dots,~$h(\vec x, y)$-এর সাহায্যে
দেওয়া, যেমন নিচের সংজ্ঞায়:
\begin{align*}
h(\vec x, 0) & = f(\vec x) \\
h(\vec x, y+1) & = g(\vec x, y, \tuple{h(\vec x, 0), \dots, h(\vec x, y)}).
\end{align*}
নিচের ছকটি ধারণাটিকে আরও সংক্ষেপে প্রকাশ করে:
\[
h(\vec x, y) = g(\vec x, y, \tuple{h(\vec x, 0), \dots, h(\vec x, y-1)})
\]
এখানে বোঝা হবে যে $y$, $0$ হলে $g$-এর শেষ আর্গুমেন্টটি কেবল শূন্য
অনুক্রম। উভয় রূপেই ধারণাটি হলো, “উত্তরসূরি ধাপ” গণনার সময় অপেক্ষক $h$
এ পর্যন্ত গণনা করা মানগুলির সম্পূর্ণ অনুক্রম ব্যবহার করতে পারে। একে
\emph{পূর্বমান-ভিত্তিক পুনরাবৃত্তি} বলা হয়। একটি নির্দিষ্ট উদাহরণে এর
সাহায্যে নিচের ধরনের সংজ্ঞাকে ন্যায্য প্রতিপন্ন করা যায়:
\begin{align*}
h(\vec x, y) & = \begin{cases}
  g(\vec x, y, h(\vec x, k(\vec x, y))) & \text{যদি $k(\vec x, y) < y$} \\
  f(\vec x) & \text{অন্যথায়।}
\end{cases}
\end{align*}
অন্যভাবে বললে, $y$-এ $h$-এর মানকে~$k$ দিয়ে প্রদত্ত
\emph{যেকোনো} পূর্ববর্তী আর্গুমেন্টে $h$-এর মানের সাহায্যে গণনা করা যায়।


\begin{prob}
  পূর্বমান-ভিত্তিক পুনরাবৃত্তির সাহায্যে ভাগশেষ অপেক্ষক $r(x,y)$
  সংজ্ঞায়িত করুন। ($x$, $y$ স্বাভাবিক সংখ্যা এবং $y > 0$ হলে $r(x,y)$
  হলো~$y$-এর চেয়ে ছোট সেই সংখ্যা, যার জন্য কোনো~$z$-এর ক্ষেত্রে
  $x = z\times y + r(x,y)$। নির্দিষ্টতার জন্য ধরা যাক, $y=0$ হলে
  $r(x,0) = 0$।)
\end{prob}

সাধারণ আদিম পুনরাবৃত্তির সাহায্যে এই অপেক্ষকগুলি কীভাবে পাওয়া যায়, তা
ভেবে দেখা উচিত। আদিম পুনরাবৃত্তির একটি শেষ রূপ আরও নমনীয়, কারণ এতে
পথ চলতে চলতে \emph{পরামিতি} (পার্শ্ব-মান) পরিবর্তন করার অনুমতি থাকে:
\begin{align*}
h(\vec x, 0) & = f(\vec x) \\
h(\vec x, y+1) & = g(\vec x, y, h(k(\vec x), y))
\end{align*}
এটিকেও সাধারণ আদিম পুনরাবৃত্তির সাহায্যে অনুকরণ করা যায়। (কাজটি
জটিল। ইঙ্গিত হিসেবে গণনাটি হাতে ধাপে ধাপে উন্মোচন করে দেখুন।)

\end{document}
